% ________________________________________________________________________________ OVERRIDEN MACROS (COMMANDS) ____________________________________________________________________________


\begin{comment}     % --- RIEPILOGO MACRO (COMANDI) MODIFICATE O SOVRASCRITTE NELLA CLASSE "HKNdocument" ---

    RIEPILOGO COMANDI SOVRASCRITTI (RENEWED) DALLA CLASSE HKNdocument
    La custom class HKNdocument modifica alcun comandi standard di LaTeX, in precisione quelli elencati a seguire, per adattarli al design una volta per tutte.
    Un comando si modifica con \renewcommand{cmd}{def}, che non fa altro che riscriverne la macro.

    --- METADATI (Titolo e Autori) ---
    \title{...}        % Invece di limitarsi a dichiarare il titolo, ora lo memorizza in una variabile interna per poterlo stampare con stili diversi nella copertina personalizzata.
    \author{...}       % Modificato per supportare più autori in sequenza: se richiami \author{} più volte, li concatena automaticamente inserendo una virgola tra l'uno e l'altro.

    --- COPERTINA E INTESTAZIONI ---
    \maketitle         % Stravolge completamente l'estetica base di LaTeX: annulla la copertina standard per generare il frontespizio ufficiale con loghi, revisori, versione e data.
    \headrule          % (Comando del pacchetto fancyhdr) Sovrascrive la riga superiore dell'intestazione, spezzandola in due segmenti separati per lasciare spazio al logo centrale.

    --- GESTIONE DEI CAPITOLI E DELLE PAGINE BIANCHE ---
    \clearpage         % Sovrascrive il salto pagina base: impone la regola \thispagestyle{empty} per assicurarsi che un'eventuale pagina bianca generata non abbia numeri o intestazioni residue.
    \cleardoublepage   % Fa la stessa cosa di \clearpage, ma per l'impaginazione fronte-retro (evitando pagine sinistre vuote ma con l'intestazione stampata).
    \chaptermark{...}  % Modifica il modo in cui il motore registra i nomi dei capitoli in memoria, permettendo all'intestazione di pagina di catturarli correttamente.

    --- STILE TIPOGRAFICO ---
    \ULdepth           % (Variabile del pacchetto ulem) Sostituisce la distanza standard della sottolineatura, abbassandola a 1.8pt per far funzionare bene il comando \myuline.

    --- NUMERAZIONE AUTOMATICA (Modificata in _preamble.tex) ---
    Numerazione Figure/Tabelle/Equazioni    % Anche se non usa renewcommand, sovrascrive la logica di base legandola ai capitoli (es. stamperà "Figura 1.1" invece di una numerazione globale come "Figura 15").

    Più in dettaglio: la numerazione standard di particolari elementi in LaTeX (immagini, tabelle, equazioni, ed eventuali altri) sfrutta di base una logica globale,
    nel senso che vengono contati da 1 ad X (es "fig. 62", presente nel capitolo 5, e non "fig 5.6", cioè la sesta figura del capitolo 5), senza che vi sia una
    gerarchia, o in precisione sistema di subordinazione, che li suddivida. Ciò è ovviamente adatto solo per documenti brevi, la cui navigazione è nativamente facile.
    Nel caso della document class "book", di cui HKNdocument è una modifica, cioè per documenti lunghi e suddivisi in capitoli, la prassi comune è quella di
    gerarchizzare i contatori in modo che siano subordinati a opportuni elementi del documento, che possono essere: part, chapter, section, subsection, sub-subsection.

    Per gli elementi standard di LaTeX questo si ottiene con il comando: \numberwithin{son}{father}, fornito dal package amsmath, che lega il contatore figlio al
    contatore padre. In precisione, in _preamble.tex caricato dalla custoclass HKNdocument, sono esplicitamente riportati i seguenti comandi per farsi che
    la numerazione di figure, tabelle, ed equazioni sia subordinata a quella dei capitoli:

    \numberwithin{figure}{chapter}
    \numberwithin{table}{chapter}
    \numberwithin{equation}{chapter}

    La particolarità di questo comando è che si tratta di una macro che al suo interno utilizza un \renewcommand, e per questo è stata riportata nella spiegazione
    relativa alle "MACRO MODIFICATE DA HKNdocument". Più precisamente questo comando sovrascrive "di nascosto" le macro originali (\thefigure, \thetable, \theequation),
    facendo stampare al posto della numerazione globale, quella subordinata al capitolo.

    Quelle appena citate sono le macro che si possono usare per "leggere" i "numeri correnti" (gli ultimi oggetti prima del comando) dei contatori. Le più utili sono:
    \thepart            % Numero della parte (solitamente numeri romani).
    \thechapter         % Numero del capitolo (non presente in article).
    \thesection         % Numero della sezione.
    \thesubsection      % Numero della sottosezione.
    \thesubsubsection   % Numero della sotto-sottosezione.
    \theparagraph       % Numero del paragrafo (spesso non numerato di default: la numerazione va attivata con \setcounter{secnumdepth}{4}, e per farli apparire anche nell'indice con \setcounter{tocdepth}{4}).
    \thesubparagraph    % Numero del sottoparagrafo (spesso non numerato di default: la numerazione va attivata con \setcounter{secnumdepth}{5}, e per farli apparire anche nell'indice con \setcounter{tocdepth}{5}).
    \thepage            % Il numero della pagina corrente.
    \theequation        % Il numero dell'equazione nell'ambiente equation.
    \thefigure          % Il numero della figura nell'ambiente figure.
    \thetable           % Il numero della tabella nell'ambiente table.
    \thefootnote        % Il simbolo o numero della nota a piè di pagina.

    Un'ultima spiegazione nota relativamente alla numerazione degli ELEMENTI CUSTOM della custom class HKNdocument, cioè dei BOX COLORATI  per Teoremi, Definizioni, Esercizi ecc.

    Per gli ambienti generati con \DeclareHKNBox, la gerarchia NON è gestita da \numberwithin. Questi box sono creati con il pacchetto 'tcolorbox', che
    possiede un suo motore indipendente. La classe gli passa semplicemente l'opzione 'number within=chapter' in fase di creazione. Il risultato estetico sul PDF è
    identico, ma il meccanismo software è nativo del pacchetto.

    RIASSUNTO FINALE COSA VIENE CAMBIATO (OVERRIDE) IN PRATICA DALLA CUSTOM CLASS HKNdocument RELATIVAMENTE ALLA NUMERAZIONE DEGLI OGGETTI:
    1. Legame gerarchico: I contatori di Figure, Tabelle ed Equazioni eccetera vengono agganciati al contatore dei Capitoli.
    2. Reset automatico: All'inizio di ogni nuovo capitolo, il conteggio degli elementi riparte da zero.
    In conclusione: Invece di usare una numerazione globale e progressiva per tutto il libro (es. "Figura 45"), il motore stamperà il numero del capitolo seguito da
    quello dell'elemento (es. "Figura 5.3", che indica a terza figura del quinto capitolo).

\end{comment}
